\subsection{Forward Chaining (DPLL-based SAT Solver)}

% Approach & Architecture
\subsubsection{Idea}
The forward-chaining solver developed by our team operates as a hybrid of
propositional Unit Propagation (the grounded counterpart of FOL Forward
Chaining) and DPLL backtracking search.

\begin{itemize}
    \item \textbf{KB construction (\textit{FOLKB}):} Converts a Futoshiki
    instance into a ground CNF Knowledge Base by fully instantiating all FOL
    axioms over the finite domain $\{1,\ldots,N\}$.  Every cell-value pair
    $(i,j,v)$ is mapped to a unique Boolean variable
    $\mathrm{var}(i,j,v) = i \cdot N^2 + j \cdot N + v \in \{1,\ldots,N^3\}$.

    \item \textbf{Forward-chaining engine (\textit{\_propagate}):} Implements
    Unit Propagation over the ground Horn KB.  Starting from the unit clauses
    that encode pre-filled cells and boundary inequalities, the engine
    repeatedly applies Modus Ponens: whenever all but one literal in a clause
    are falsified, the remaining literal is forced \emph{True}.  This is
    exactly FOL Forward Chaining restricted to the grounded domain.

    \item \textbf{Search controller (\textit{\_dpll}):} When Unit Propagation
    can no longer derive new facts, the controller selects an unassigned
    variable via the MRV\,+\,Frequency heuristic, branches on its two possible
    truth values, and backtracks upon conflict.
\end{itemize}

\paragraph{Why Unit Propagation equals Forward Chaining on a ground KB:}
The fundamental FOL axiom ``every cell has at least one value''
($\forall i\,\forall j\,\exists v\;Val(i,j,v)$) contains an existential
quantifier that cannot be expressed as a definite Horn rule in its general
form.  After \emph{full grounding} over the finite domain $\{1,\ldots,N\}$,
however, this axiom becomes the finite disjunction
$\bigvee_{v=1}^{N} x_{i,j,v}$, which \emph{is} a valid CNF clause.
Consequently, every step of Unit Propagation on the ground KB is logically
equivalent to one application of Modus Ponens on the corresponding grounded
Horn rule --- the defining operation of Forward Chaining.

\paragraph{Technical Implementation Notes:}~
\begin{itemize}
    \item \textbf{0-based Indexing:} While the formal FOL axioms in this
    report use 1-based indexing (as per the mathematical specification
    $1..N$), the internal variable encoding uses 0-based row/column indices.
    This aligns naturally with Python's list structures and eliminates
    redundant $\pm1$ arithmetic at every node expansion.

    \item \textbf{Full Grounding vs.\ Partial Grounding:} Unlike the Backward
    Chaining solver which keeps some rules in first-order form, this solver
    grounds the entire KB upfront.  For grid sizes $N \le 9$ the resulting CNF
    (at most $N^3 \approx 729$ variables and $O(N^4)$ clauses) remains
    tractable and enables $O(1)$ variable lookup via a direct-indexed model
    array.

    \item \textbf{Literal-to-clause index:} An inverted index
    $\mathrm{literal\_to\_clauses}[\ell]$ stores exactly those clauses that
    contain literal $\ell$.  When literal $p$ is set \emph{True}, only the
    clauses in $\mathrm{literal\_to\_clauses}[-p]$ need to be re-examined,
    reducing per-step cost from $O(|\mathcal{C}|)$ to $O(\deg(p))$.

    \item \textbf{Step Callback for visualisation:}
    $\mathrm{step\_callback}(r,c,v)$ is invoked each time a cell is assigned
    ($v > 0$) or retracted ($v = 0$), enabling real-time GUI updates without
    coupling the solver to the display layer.
\end{itemize}

% -----------------------------------------------------------------------
\subsubsection{Algorithm Analysis}

\paragraph{CNF Encoding --- Variable Schema}~\\
Each Boolean variable encodes one cell-value pair:
\[
    \mathrm{var}(i,\,j,\,v) \;=\; i \cdot N^2 + j \cdot N + v,
    \qquad i,j \in \{0,\ldots,N-1\},\quad v \in \{1,\ldots,N\}.
\]
The model array $\mathit{model}[1..N^3]$ stores $\{-1,\,0,\,+1\}$ for
\emph{False}, \emph{unassigned}, and \emph{True} respectively.

\paragraph{CNF Clause Groups}~\\
The ground KB is built from nine constraint groups:
\begin{center}
\begin{tabular}{lll}
\hline
\textbf{Group} & \textbf{Clause form} & \textbf{Constraint} \\
\hline
A1 & $\bigl(x_{i,j,1} \lor \cdots \lor x_{i,j,N}\bigr)$ & At-least-one per cell \\
A2 & $(\neg x_{i,j,a} \lor \neg x_{i,j,b}),\; a<b$ & At-most-one per cell \\
A3 & $(\neg x_{i,j_1,v} \lor \neg x_{i,j_2,v}),\; j_1<j_2$ & Row uniqueness \\
A4 & $(\neg x_{i_1,j,v} \lor \neg x_{i_2,j,v}),\; i_1<i_2$ & Column uniqueness \\
A5/A6 & $(\neg x_{i,j,a} \lor \neg x_{i,j+1,b})$ if $a \not< b$ & Horizontal $<$ / $>$ \\
A7/A8 & $(\neg x_{i,j,a} \lor \neg x_{i+1,j,b})$ if $a \not< b$ & Vertical $\wedge$ / $\vee$ \\
A9 & $(x_{i,j,v})$ & Pre-filled clues \\
A10 & \emph{(implicit in encoding)} & Domain $[1..N]$ \\
A11 & \emph{(implicit, via negation)} & Inequality via $\neg(a < b)$ \\
\hline
\end{tabular}
\end{center}

\paragraph{Unit Propagation (Forward Chaining step)}~\\
When literal $p$ is assigned \emph{True}, every clause in
$\mathrm{literal\_to\_clauses}[-p]$ is inspected:
\begin{itemize}
    \item If the clause already contains a \emph{True} literal: skip.
    \item If exactly one literal $\ell$ remains unassigned: add $\ell$ to the
    agenda (Modus Ponens / unit propagation step).
    \item If no unassigned literal remains and the clause is unsatisfied:
    \textbf{conflict} --- return \emph{False}.
\end{itemize}

\paragraph{Variable Selection Heuristic: MRV + Frequency}~\\
When Unit Propagation halts without completing the assignment, the solver
must branch.  Variable selection proceeds in two levels:
\begin{enumerate}
    \item \textbf{Minimum Remaining Values (MRV):} Among all unassigned cells,
    choose the cell $(i,j)$ whose \emph{domain size}
    $|\{v : \mathit{model}[\mathrm{var}(i,j,v)] \neq -1\}|$ is smallest.
    This \emph{fail-first} principle detects dead ends as early as possible.
    \item \textbf{Frequency tiebreaker:} Among all domain literals of the
    chosen cell, select the literal with the highest frequency in the CNF ---
    assigning a frequent literal satisfies more clauses simultaneously.
\end{enumerate}

\paragraph{Complexity Analysis}~
\begin{center}
\begin{tabular}{llp{6cm}}
\hline
\textbf{Aspect} & \textbf{Bound} & \textbf{Remark} \\
\hline
Variables & $N^3$ & One Boolean per cell-value pair \\
Clauses (A1) & $N^2$ & One per cell \\
Clauses (A2) & $N^2 \binom{N}{2}$ & Pairwise exclusion per cell \\
Clauses (A3+A4) & $2N\binom{N}{2}N$ & Row and column uniqueness \\
Max recursion depth & $\le N^2$ & One cell resolved per level \\
Time (worst case) & $O(N^3 \cdot 2^{N^3})$ & Full Boolean search space \\
Time (average) & $O(N^3 \cdot C)$ & $C$ = actual backtracks (small) \\
Space & $O(N^3)$ & Model array + agenda \\
\hline
\end{tabular}
\end{center}

\paragraph{Advantages and Limitations}~
\begin{itemize}
    \item \textbf{Complete:} Backtracking guarantees a solution is found if
    one exists; the solver never misses a valid assignment.
    \item \textbf{Efficient propagation:} The inverted index restricts each
    propagation step to $O(\deg(p))$ clause checks instead of $O(|\mathcal{C}|)$.
    \item \textbf{Low memory footprint:} Only the flat model array and a small
    agenda deque are maintained; no explicit clause deletion is performed.
    \item \textbf{No clause learning (CDCL):} The solver does not learn
    \emph{conflict clauses}, unlike modern industrial SAT solvers (MiniSat,
    Glucose).  Repeated conflicts in similar sub-trees are not avoided.
    \item \textbf{Static index:} $\mathrm{literal\_to\_clauses}$ is built once
    and not updated; satisfied clauses are still visited during propagation,
    incurring a minor constant overhead.
\end{itemize}

% -----------------------------------------------------------------------
\subsubsection{Pseudocode}

% SOLVE
\begin{algorithm}[H]
\caption{Master Coordinator}
\label{alg:fc-solve}
\begin{algorithmic}[1]
\Function{solve}{\textit{puzzle}}
    \State $kb \gets \Call{FOLKB}{puzzle}$
    \State $\mathit{clauses} \gets kb.\Call{build\_KB}{}$
    \State $\mathit{max\_var} \gets N^3$
    \State $\mathit{lit\_to\_clauses} \gets \Call{BuildIndex}{\mathit{clauses}}$
    \State $\mathit{freq} \gets \Call{CountFrequency}{\mathit{clauses}}$
    \State $\mathit{model} \gets \mathbf{array}[0 \ldots \mathit{max\_var}]$ initialised to $0$
    \State $\mathit{agenda} \gets \Call{deque}{}$
    \For{each clause $c \in \mathit{clauses}$ with $|c| = 1$}
        \State $\mathit{agenda}.\Call{append}{c[0]}$
    \EndFor
    \State $\mathit{result} \gets \Call{DPLL}{\mathit{clauses},\,\mathit{lit\_to\_clauses},\,\mathit{agenda},\,\mathit{model}}$
    \If{$\mathit{result} \neq \mathit{None}$}
        \State \Return $\Call{TranslateToGrid}{\mathit{result}}$
    \Else
        \State \Return $\mathit{None}$ \Comment{Puzzle has no solution}
    \EndIf
\EndFunction
\end{algorithmic}
\end{algorithm}

% DPLL
\begin{algorithm}[H]
\caption{DPLL --- Recursive Branch-and-Propagate}
\label{alg:fc-dpll}
\begin{algorithmic}[1]
\Function{DPLL}{\textit{clauses}, \textit{index}, \textit{agenda}, \textit{model}}
    \State $\mathit{nodes\_expanded} \mathrel{+}= 1$
    \If{\textbf{not} $\Call{Propagate}{\mathit{clauses},\,\mathit{index},\,\mathit{agenda},\,\mathit{model}}$}
        \State \Return $\mathit{None}$ \Comment{Conflict --- backtrack}
    \EndIf
    \State $\mathit{best\_var} \gets \Call{GetUnassignedVar}{\mathit{model}}$
    \If{$\mathit{best\_var} = \mathit{None}$}
        \State \Return $\mathit{model}$ \Comment{All variables assigned --- solution found}
    \EndIf
    % Branch True
    \State $\mathit{agenda\_true} \gets \Call{deque}{[\mathit{best\_var}]}$
    \State $\mathit{result} \gets \Call{DPLL}{\mathit{clauses},\,\mathit{index},\,\mathit{agenda\_true},\,\mathit{model}}$
    \If{$\mathit{result} \neq \mathit{None}$}
        \State \Return $\mathit{result}$
    \EndIf
    \State \Comment{Branch False --- restore snapshot}
    \State $\mathit{model} \gets \mathit{saved}$
    \State $\mathit{agenda\_false} \gets \Call{deque}{[-\mathit{best\_var}]}$
    \State \Return $\Call{DPLL}{\mathit{clauses},\,\mathit{index},\,\mathit{agenda\_false},\,\mathit{model}}$
\EndFunction
\end{algorithmic}
\end{algorithm}

% PROPAGATE
\begin{algorithm}[H]
\caption{Unit Propagation --- Forward Chaining Engine}
\label{alg:fc-propagate}
\begin{algorithmic}[1]
\Function{Propagate}{\textit{clauses}, \textit{index}, \textit{agenda}, \textit{model}}
    \While{$\mathit{agenda}$ is not empty}
        \State $p \gets \mathit{agenda}.\Call{popleft}{}$
        \If{$\Call{IsFalse}{p,\,\mathit{model}}$}
            \State \Return $\mathit{False}$ \Comment{Immediate conflict}
        \EndIf
        \If{\textbf{not} $\Call{IsAssigned}{|p|,\,\mathit{model}}$}
            \State $\Call{Assign}{p,\,\mathit{model}}$ \Comment{Modus Ponens: set literal True}
            \If{callback $\neq$ None \textbf{and} $p > 0$}
                \State $\Call{callback}{\mathrm{Decode}(p)}$ \Comment{Notify GUI}
            \EndIf
        \EndIf
        \For{each clause $c \in \mathit{index}[-p]$}
            \State $\mathit{sat} \gets \mathit{False}$;
                   $\;\mathit{unassigned} \gets 0$;
                   $\;\mathit{unit} \gets 0$
            \For{each literal $\ell \in c$}
                \If{$\Call{IsTrue}{\ell,\,\mathit{model}}$}
                    \State $\mathit{sat} \gets \mathit{True}$; \textbf{break}
                \EndIf
                \If{\textbf{not} $\Call{IsAssigned}{|\ell|,\,\mathit{model}}$}
                    \State $\mathit{unassigned} \mathrel{+}= 1$;
                           $\;\mathit{unit} \gets \ell$
                    \If{$\mathit{unassigned} > 1$} \textbf{break} \EndIf
                \EndIf
            \EndFor
            \If{$\mathit{sat}$} \textbf{continue} \EndIf
            \If{$\mathit{unassigned} = 0$} \Return $\mathit{False}$ \Comment{CONFLICT} \EndIf
            \If{$\mathit{unassigned} = 1$ \textbf{and} $\mathit{unit} \notin \mathit{agenda\_set}$}
                \State $\mathit{agenda}.\Call{append}{\mathit{unit}}$ \Comment{New unit clause}
            \EndIf
        \EndFor
    \EndWhile
    \State \Return $\mathit{True}$
\EndFunction
\end{algorithmic}
\end{algorithm}

% GET UNASSIGNED VAR
\begin{algorithm}[H]
\caption{Variable Selection --- MRV + Frequency Heuristic}
\label{alg:fc-heuristic}
\begin{algorithmic}[1]
\Function{GetUnassignedVar}{\textit{model}}
    \State $\mathit{best\_lit} \gets \mathit{None}$;
           $\;\mathit{best\_dom} \gets \infty$;
           $\;\mathit{best\_freq} \gets -1$
    \For{$i \in 0 \ldots N-1$}
        \For{$j \in 0 \ldots N-1$}
            \If{$\exists\,v : \mathit{model}[\mathrm{var}(i,j,v)] = 1$}
                \textbf{continue} \Comment{Cell already assigned}
            \EndIf
            \State $\mathit{dom} \gets [\mathrm{var}(i,j,v) \mid v \in 1..N,\;
                   \mathit{model}[\mathrm{var}(i,j,v)] \neq -1]$
            \If{$\mathit{dom} = \emptyset$} \textbf{continue} \EndIf
            \State $\mathit{cand} \gets \arg\max_{\ell \in \mathit{dom}}\;\mathit{freq}[\ell]$
            \State $f \gets \mathit{freq}[\mathit{cand}]$
            \If{$|\mathit{dom}| < \mathit{best\_dom}$ \textbf{or}
                $(|\mathit{dom}| = \mathit{best\_dom}$ \textbf{and} $f > \mathit{best\_freq})$}
                \State $\mathit{best\_dom} \gets |\mathit{dom}|$;
                       $\;\mathit{best\_lit} \gets \mathit{cand}$;
                       $\;\mathit{best\_freq} \gets f$
            \EndIf
        \EndFor
    \EndFor
    \State \Return $\mathit{best\_lit}$ \Comment{$\mathit{None}$ if all cells are assigned}
\EndFunction
\end{algorithmic}
\end{algorithm}

% BUILD INDEX + COUNT FREQ
\begin{algorithm}[H]
\caption{Index Construction and Frequency Count}
\label{alg:fc-index}
\begin{algorithmic}[1]
\Function{BuildIndex}{\textit{clauses}}
    \State $\mathit{index} \gets \mathbf{defaultdict}(\mathit{list})$
    \For{each clause $c \in \mathit{clauses}$}
        \For{each literal $\ell \in c$}
            \State $\mathit{index}[\ell].\Call{append}{c}$
        \EndFor
    \EndFor
    \State \Return $\mathit{index}$
\EndFunction

\Function{CountFrequency}{\textit{clauses}}
    \State $\mathit{freq} \gets \mathbf{defaultdict}(0)$
    \For{each clause $c \in \mathit{clauses}$}
        \For{each literal $\ell \in c$}
            \State $\mathit{freq}[|\ell|] \mathrel{+}= 1$
        \EndFor
    \EndFor
    \State \Return $\mathit{freq}$
\EndFunction

\Function{TranslateToGrid}{\textit{model}}
    \State $\mathit{grid} \gets N \times N$ matrix of $0$
    \For{$i \in 0..N-1$}
        \For{$j \in 0..N-1$}
            \For{$v \in 1..N$}
                \If{$\mathit{model}[\mathrm{var}(i,j,v)] = 1$}
                    \State $\mathit{grid}[i][j] \gets v$; \textbf{break}
                \EndIf
            \EndFor
        \EndFor
    \EndFor
    \State \Return $\mathit{grid}$
\EndFunction
\end{algorithmic}
\end{algorithm}

\subsubsection{Illustration}
[Chèn nội dung tại đây]
